\documentclass[11pt]{article}

\usepackage{amsmath,amssymb,bm}
\usepackage{geometry}
\usepackage{physics}
\usepackage{amsfonts}

\geometry{margin=1in}

\title{Implicit DG/HDG Algorithm for the 2D SWE with Wetting and Drying}
\author{}
\date{}

\begin{document}

\maketitle

You’re right: Exasim’s DG/HDG time integration is typically implicit (Backward Euler / DIRK + Newton–Krylov), so the ``explicit-update → limiter → dry-enforcement'' workflow needs to be adapted.

A robust implicit wetting/drying (W/D) scheme for 2D SWE hinges on one key design choice:

How you enforce $H \ge 0$ inside Newton iterations.

Below are two practical, working approaches used in implicit solvers. I present an implicit HDG algorithm (with static condensation) and show where wet/dry logic enters.

\section*{0) Problem setup}

Unknowns (conservative):
\[
\mathbf{U} = (H,\ q_x,\ q_y), \qquad \mathbf{q} = H\mathbf{u}.
\]

Let the semi-discrete HDG residual be
\[
\mathcal{R}(\mathbf{U}, \widehat{\mathbf{U}};\ z_b,\text{params}) = 0,
\]
where $\widehat{\mathbf{U}}$ are trace unknowns on faces (global), and $\mathbf{U}$ are element unknowns (local).

Implicit time step (Backward Euler shown; DIRK is similar stage-by-stage):
\[
\mathbf{G}(\mathbf{U}^{n+1}, \widehat{\mathbf{U}}^{n+1})
=
\frac{\mathbf{U}^{n+1}-\mathbf{U}^n}{\Delta t}
+\mathcal{R}(\mathbf{U}^{n+1},\widehat{\mathbf{U}}^{n+1})
=0.
\]

\section*{1) The central difficulty in implicit W/D}

In Newton iterations, intermediate iterates can produce:
\begin{itemize}
\item negative $H$,
\item huge $\mathbf{u}=\mathbf{q}/H$,
\item invalid wave speeds $\sqrt{gH}$,
\item divergence of the nonlinear solve.
\end{itemize}

So for implicit methods, you typically do one of:

\subsection*{Option A (recommended): Depth transform to guarantee positivity}

Represent depth as
\[
H = H_{\min} + \exp(\phi)
\]
(or softplus), and solve for $\phi$ instead of $H$.

This makes $H>H_{\min}$ automatically for every Newton iterate, which is extremely stabilizing.

\subsection*{Option B: Active-set / projection (semi-smooth Newton)}

Treat dryness as an inequality constraint:
\[
H \ge 0,\quad \mathbf{q}=0\ \text{when dry},
\]
and update the wet/dry set inside the nonlinear iterations.

This can be more ``physically sharp'' but is harder to implement robustly.

Given you want something robust in Exasim, Option A is usually the best starting point.

\section*{2) Implicit HDG algorithm with positivity via depth transform (Option A)}

\subsection*{Unknowns}

Solve for:
\[
\mathbf{W}=(\phi,\ q_x,\ q_y),\quad \widehat{\mathbf{W}}=(\widehat{\phi},\ \widehat{q}_x,\ \widehat{q}_y),
\]
and map to conservative variables:
\[
H(\phi)=H_{\min}+\exp(\phi),\qquad \mathbf{U}(\mathbf{W})=(H(\phi),q_x,q_y).
\]

\subsection*{Wet/dry classification (used only for physics switches)}

Define a smooth wet indicator to avoid discontinuities in Newton:
\[
\chi(H)=\frac{1}{2}\left(1+\tanh\left(\frac{H-H_{\text{dry}}}{\varepsilon}\right)\right)\in(0,1),
\]
with $\varepsilon$ small (e.g. $0.1H_{\text{dry}}$).

Use $\chi$ to smoothly turn on/off terms that are problematic near dry states (e.g., friction, wind stress, certain closures).

\subsection*{Regularized velocity used in fluxes/closures}

\[
\mathbf{u} = \frac{\mathbf{q}}{H_{\text{reg}}},\qquad
H_{\text{reg}}=\max(H, H_{\text{dry}})
\]
(or a smooth max).

\subsection*{Hydrostatic reconstruction in fluxes (well-balanced + W/D safe)}

On each face quadrature point, use reconstructed nonnegative depths built from $\eta$ and bed elevation (your convention consistent):
\[
\eta = H + z_b,\qquad z_b^\*=\max(z_{b,L},z_{b,R}),
\]
\[
H_L^\*=\max(H_{\min},\eta_L-z_b^\*),\qquad
H_R^\*=\max(H_{\min},\eta_R-z_b^\*),
\]
and reconstructed momenta $\mathbf{q}^\* = H^\*\mathbf{u}$ with regularized $\mathbf{u}$.

\subsection*{Implicit step: Newton–Krylov with HDG condensation}

At time step $n\to n+1$, solve:
\[
\mathbf{F}(\mathbf{W},\widehat{\mathbf{W}})
=
\frac{\mathbf{U}(\mathbf{W})-\mathbf{U}^n}{\Delta t}
+
\mathcal{R}(\mathbf{U}(\mathbf{W}), \widehat{\mathbf{U}}(\widehat{\mathbf{W}}))
=0.
\]

\subsubsection*{Algorithm block (Backward Euler; HDG)}

\begin{enumerate}
\item \textbf{Initialize} $\mathbf{W}^{(0)}$ from $\mathbf{U}^n$:  
$\phi^{(0)}=\log(\max(H^n-H_{\min},\text{tiny}))$, $\mathbf{q}^{(0)}=\mathbf{q}^n$.  
Initialize $\widehat{\mathbf{W}}^{(0)}$ similarly (or extrapolate).

\item \textbf{Newton loop} for $k=0,1,\dots$:
\begin{itemize}
\item Local element solves (static condensation preparation):  
For each element $K$, given trace guess $\widehat{\mathbf{W}}^{(k)}$, compute the element residual and local Jacobian blocks
\[
\mathbf{F}_K(\mathbf{W}_K,\widehat{\mathbf{W}}_K)=0,
\]
including:
\begin{itemize}
\item hydrostatic reconstruction in numerical flux,
\item regularized velocity $\mathbf{u}(\mathbf{W})$,
\item smooth wet indicator $\chi(H)$ gating stiff/degenerate terms.
\end{itemize}

\item Condense element unknown increments $\delta \mathbf{W}_K$ to obtain a global system for trace increments $\delta \widehat{\mathbf{W}}$:
\[
\mathbf{S}\,\delta \widehat{\mathbf{W}} = -\mathbf{g},
\]
where $\mathbf{S}$ is the HDG Schur complement.

\item Solve the global linear system with GMRES (preconditioned).

\item Recover local increments $\delta \mathbf{W}_K$ element-by-element.

\item Line search / damping to ensure stability of Newton:
\[
\mathbf{W}^{(k+1)}=\mathbf{W}^{(k)}+\alpha \delta \mathbf{W},\quad
\widehat{\mathbf{W}}^{(k+1)}=\widehat{\mathbf{W}}^{(k)}+\alpha \delta \widehat{\mathbf{W}},
\]
with $\alpha\in(0,1]$ chosen to reduce $\|\mathbf{F}\|$.  
(Positivity is automatic because $H=H_{\min}+\exp(\phi)$.)

\item Converge when $\|\mathbf{F}\|$ and/or $\|\delta\|$ are below tolerance.
\end{itemize}

\item \textbf{Post-step ``hard'' dryness option (optional)}:  
If you want strictly dry cells (not thin film), you can apply:
if $H < H_{\text{dry}}$ then set $\mathbf{q}=0$ (and optionally clamp $H\to H_{\min}$) after convergence.  
This is not required, but sometimes helpful for clean dry masks.
\end{enumerate}

Why this works well in implicit HDG:
\begin{itemize}
\item Newton never sees negative $H$.
\item Flux wave speeds remain real.
\item The only non-smoothness left is in wet/dry switching, which we smooth with $\chi(H)$.
\end{itemize}

\section*{3) Implicit HDG algorithm with an active-set dry mask (Option B)}

If you insist on true dry states ($H=0$) inside the solve (no thin film), you can do an active-set strategy.

\subsection*{Key idea}

At Newton iteration $k$, define wet set:
\[
\mathcal{W}^{(k)}=\{K:\ \bar H_K^{(k)} > H_{\text{dry}}\},\qquad
\mathcal{D}^{(k)}=\Omega\setminus\mathcal{W}^{(k)}.
\]

Then enforce in dry elements:
\[
H=0,\quad \mathbf{q}=0,
\]
and remove those DOFs from the Newton update (or strongly penalize them).

\subsection*{Semi-smooth Newton / active set loop}

At each Newton iteration:
\begin{enumerate}
\item Compute $\bar H_K^{(k)}$, update $\mathcal{W}^{(k)},\mathcal{D}^{(k)}$.
\item Assemble residual/Jacobian:
\begin{itemize}
\item In wet cells: full SWE residual.
\item In dry cells: constraints residual $H=0,\mathbf{q}=0$.
\end{itemize}
\item Solve condensed HDG linear system for increments.
\item Apply line search that prevents ``unphysical'' creation of negative depths.
\item Repeat until both:
\begin{itemize}
\item residual small,
\item active set no longer changes.
\end{itemize}
\end{enumerate}

This is more complex and less forgiving than the transform approach, but it can give a sharper shoreline.

\section*{4) What I recommend for Exasim}

Start with Option A (depth transform) plus:
\begin{itemize}
\item hydrostatic reconstruction in fluxes,
\item smooth wet indicator $\chi(H)$,
\item regularized division $\mathbf{q}/H_{\text{reg}}$,
\item implicit treatment of friction (it fits naturally in BE/IRK).
\end{itemize}

This combination is the most robust path to a convergent implicit solver with W/D.

\section*{5) Minimal implicit HDG W/D algorithm block you can implement first}

\begin{enumerate}
\item Use Backward Euler.
\item Solve with Newton–GMRES + HDG condensation.
\item Enforce $H>0$ via $H=H_{\min}+\exp(\phi)$.
\item Use hydrostatic reconstruction for face fluxes.
\item Gate stiff terms with $\chi(H)$.
\end{enumerate}

That gets you a stable, production-capable baseline.

If you tell me which Exasim SWE form you’re using (e.g., $(H,Hu,Hv)$ with bathymetry source split, and your exact bed convention $z_b$ vs $h$), I can write the exact HDG residual form $\mathcal{R}$ and its block Jacobian structure suitable for static condensation (including which terms to lag vs differentiate for better Newton robustness).


\

\

\noindent 
{\Large \textbf{2D Nonlinear Shallow Water Equations in Transformed Variables $(\phi,q_x,q_y)$}}

\


Below is the 2D nonlinear SWE rewritten with the positivity-enforcing transformation
\[
H(\phi)=H_{\min}+e^{\phi},\qquad H_{\min}\ge 0,
\]
using the unknowns $(\phi,q_x,q_y)$ where $\mathbf{q}=(q_x,q_y)$.

I give (i) the conservative form with transformed depth, (ii) the expanded $\phi$-equation, and (iii) a clean vector form.

\section*{1) Start from conservative SWE in $(H,\mathbf{q})$}

Mass:
\[
\partial_t H+\nabla\cdot\mathbf{q}=0.
\]

Momentum:
\[
\partial_t \mathbf{q}+\nabla\cdot\!\left(\frac{\mathbf{q}\otimes\mathbf{q}}{H}+\frac{1}{2}gH^2\mathbf{I}\right)
=-gH\nabla z_b+\mathbf{S}_{phys},
\]
where $z_b(x,y)$ is bed elevation (positive upward) and $\eta=H+z_b$.

\section*{2) Substitute $H = H_{\min}+e^{\phi}$}

Define
\[
H(\phi)=H_{\min}+e^{\phi},\qquad
\partial_t H = e^{\phi}\,\partial_t\phi,\qquad
\nabla H = e^{\phi}\,\nabla\phi.
\]

\subsection*{2.1 Continuity equation in $\phi$}

From $\partial_t H+\nabla\cdot\mathbf{q}=0$:
\[
e^{\phi}\,\partial_t\phi+\nabla\cdot\mathbf{q}=0.
\]

Equivalently,
\[
\boxed{\;\partial_t\phi=-e^{-\phi}\,\nabla\cdot\mathbf{q}\;}
\]
(note the singularity as $\phi\to -\infty$ corresponds to $H\to H_{\min}$; numerically you often keep $H_{\min}>0$ and/or regularize).

\section*{3) Momentum equations in $(\phi,q_x,q_y)$}

The momentum PDE remains an evolution equation for $\mathbf{q}$, with $H(\phi)$ inserted wherever $H$ appears.

Let
\[
H(\phi)=H_{\min}+e^{\phi}.
\]

Then
\[
\boxed{
\partial_t \mathbf{q}
+\nabla\cdot\!\left(\frac{\mathbf{q}\otimes\mathbf{q}}{H(\phi)}
+\frac{1}{2}g\,H(\phi)^2\,\mathbf{I}\right)
=
-g\,H(\phi)\,\nabla z_b
+\mathbf{S}_{phys}(\phi,\mathbf{q})
}.
\]

Written componentwise ($i=x,y$):
\[
\partial_t q_i
+
\partial_x\!\left(\frac{q_i q_x}{H(\phi)}\right)
+
\partial_y\!\left(\frac{q_i q_y}{H(\phi)}\right)
+
\partial_i\!\left(\frac{1}{2}gH(\phi)^2\right)
=
-gH(\phi)\,\partial_i z_b
+S_i.
\]

You can combine the pressure gradient and bed slope term using $\eta=H+z_b$:
\[
\partial_i\!\left(\frac{1}{2}gH^2\right)+gH\,\partial_i z_b
= gH\,\partial_i(H+z_b)=gH\,\partial_i\eta,
\]
so the momentum equation can also be expressed as
\[
\boxed{
\partial_t \mathbf{q}
+\nabla\cdot\!\left(\frac{\mathbf{q}\otimes\mathbf{q}}{H(\phi)}\right)
+g\,H(\phi)\,\nabla \eta
=
\mathbf{S}_{phys}
}.
\]
(Choose the form that best matches your well-balanced discretization; hydrostatic reconstruction is usually built around $\eta$.)

\section*{4) Full transformed system (compact)}

Let the state be
\[
\mathbf{W}=(\phi,q_x,q_y)^T,\qquad H(\phi)=H_{\min}+e^{\phi}.
\]

Then
\[
\boxed{
\begin{aligned}
e^{\phi}\,\partial_t\phi+\nabla\cdot\mathbf{q}&=0,\\[4pt]
\partial_t \mathbf{q}
+\nabla\cdot\!\left(\frac{\mathbf{q}\otimes\mathbf{q}}{H(\phi)}
+\frac{1}{2}gH(\phi)^2\mathbf{I}\right)
&=-gH(\phi)\nabla z_b+\mathbf{S}_{phys}(\phi,\mathbf{q}).
\end{aligned}
}
\]

We begin with the transformed continuity equation

\[
\partial_t\phi + e^{-\phi}\,\nabla\cdot\mathbf{q} = 0.
\]

Let 
\[
a(\phi)=e^{-\phi}.
\]

Then the vector calculus identity (product rule) states

\[
\nabla\cdot(a\,\mathbf{q}) 
= a\,\nabla\cdot\mathbf{q} + \mathbf{q}\cdot\nabla a.
\]

Rearranging,

\[
a\,\nabla\cdot\mathbf{q}
=
\nabla\cdot(a\,\mathbf{q})
-
\mathbf{q}\cdot\nabla a.
\]

Substituting $a = e^{-\phi}$ into the transformed continuity equation gives

\[
\partial_t\phi
+
\nabla\cdot(e^{-\phi}\mathbf{q})
-
\mathbf{q}\cdot\nabla(e^{-\phi})
= 0.
\]

Now compute the gradient:

\[
\nabla(e^{-\phi})
=
- e^{-\phi}\nabla\phi.
\]

Therefore,

\[
-\mathbf{q}\cdot\nabla(e^{-\phi})
=
-\mathbf{q}\cdot(-e^{-\phi}\nabla\phi)
=
e^{-\phi}\,\mathbf{q}\cdot\nabla\phi.
\]

Thus the correctly expanded form is

\[
\boxed{
\partial_t\phi
+
\nabla\cdot(e^{-\phi}\mathbf{q})
+
e^{-\phi}\,\mathbf{q}\cdot\nabla\phi
= 0.
}
\]


Then
\[
\boxed{
\begin{aligned}
\partial_t\phi
+
\nabla\cdot(e^{-\phi}\mathbf{q})
+
e^{-\phi}\,\mathbf{q}\cdot\nabla\phi
& = 0,\\[4pt]
\partial_t \mathbf{q}
+\nabla\cdot\!\left(\frac{\mathbf{q}\otimes\mathbf{q}}{H(\phi)}
+\frac{1}{2}gH(\phi)^2\mathbf{I}\right)
&=-gH(\phi)\nabla z_b+\mathbf{S}_{phys}(\phi,\mathbf{q}).
\end{aligned}
}
\]


\section*{5) Notes for implicit Newton/HDG implementation}

\begin{itemize}
\item The map $H(\phi)$ gives automatic positivity $H\ge H_{\min}$ for all Newton iterates.
\item For Jacobians you need:
\[
\frac{dH}{d\phi}=e^{\phi},\quad
\frac{d}{d\phi}\left(\frac{1}{H(\phi)}\right)=-\frac{e^{\phi}}{H(\phi)^2},\quad
\frac{d}{d\phi}\left(H(\phi)^2\right)=2H(\phi)e^{\phi}.
\]
\item If you want true dry states, you can take $H_{\min}$ very small and separately apply a ``hard dry'' postprocessing ($\mathbf{q}\to 0$ when $H<H_{\text{dry}}$), but the PDE system above remains well-defined.
\end{itemize}

If you want, I can also write the transformed system in flux form for $\mathbf{W}=(\phi,q_x,q_y)$ suitable for HDG assembly (i.e., identify the fluxes and source terms explicitly and provide the derivatives needed for Newton linearization).


\end{document}